% Page 024 — page imprimée 20
% Protestants et Catholiques face à la Révolution à Meyrueis et dans les Cévennes


{\fontsize{10.15}{11.65}\selectfont

\begin{center}
{\Large\bfseries Protestants et Catholiques face à la Révolution\\
à Meyrueis et dans les Cévennes}

\vspace{1.0em}

{\bfseries 1789 -- 1793}
\end{center}

\vspace{0.55em}

Meyrueis, comme toutes les Cévennes méridionales, s’engage dans la Révolution avec ardeur.
La majorité de la population ayant abjuré du bout des lèvres, était restée attachée à la Réforme et avait accueilli avec joie l’Edit de Tolérance de 1797 qui leur permettait de retrouver un état civil pour leurs familles.
Avec la Révolution les protestants étaient désormais des citoyens à part entière. Les campagnes cévenoles envoyèrent donc à l’Assemblée Constituante une majorité de « patriotes ».

Mais l’année 1789 avait été une mauvaise année : les récoltes médiocres après un hiver glacial et les difficultés économiques de toutes sortes provoquèrent des troubles. Uns grande peur touche aussi les communautés protestantes qui ont tout à craindre d’une réaction catholique.
Les persécutions et les dragonnades sont dans toutes les mémoires.
Les différences confessionnelles imprègnent celles des comportements. Le choc décisif est la question religieuse.

\medskip

\textbf{La réaction catholique}

\vspace{0.45em}

L’Assemblée Constituante vota, en effet, la vente des biens nationaux et la Constitution Civile du Clergé exige de tous les prêtres un serment de fidélité à la Constitution. Dès lors les prêtres se divisèrent en « constitutionnels », ou « réfractaires » nombreux dans le nord de la Lozère.
Le pape Pie VI condamnant cette constitution, il y eut schisme dans l’église catholique : la plupart des prêtres réfractaires prennent le parti de la contre-révolution : certains émigrèrent d’autres furent massacrés pendant la Terreur (1792-1794).

Comme les chouans en Vendée, le nord de la Lozère, à partir des Gorges du Tarn se dressa contre la République laïque et anticléricale. Des bandes se formèrent pour lutter contre les « patriotes ». Une véritable petite armée était disponible pour des coups de main ou de véritables batailles, dirigées par un chef noble, originaire de Nasbinals, Marc-Antoine Charrier, député du Gévaudan.

Des « républicains » furent tués et la situation devint si grave que la République envoya des troupes pour stopper « les rebelles ». Beaucoup furent arrêtés dont les principaux chefs.
Charrier s’était aménagé une « cache souterraine » dans sa maison de Nasbinals. Mais il fut vendu et il sera fusillé à Rodez. C’est ainsi que dans la région

\vspace{0.25em}

{\fontsize{8.0}{9.25}\selectfont
\textit{« Jean-Baptiste Fages de la Malène, frère d’un rebelle guillotiné à Florac en Juin 1793, est le chef de la première bande organisée.
Ses hommes se retrouvent lors d’expéditions punitives contre « les patriotes » et le reste du temps réintègrent leurs foyers et cultivent leurs champs.
Le lieutenant de Fages est un jeune déserteur des armées républicaines, natif de Hures, sur le Causse Méjean connu sous le sobriquet de Ferratou. Tous deux sont arrêtés en Janvier 1794. Fages périt à l’échafaud alors que Ferratou parvient à s’échapper. A la tête de sa bande il fut repris un des jours suivant et fusillé à Meyrueis ».}
}

}

